

\begin{figure*}[!t]
\centering
\begin{adjustbox}{max width=0.86\textwidth}
\begin{tikzpicture}[
  font=\small,
  stage/.style={draw=tealrule, line width=0.7pt, rounded corners=2pt,
                fill=tealrow, text width=4.9cm, align=center, inner sep=5pt},
  excl/.style={draw=gray!55, line width=0.6pt, rounded corners=2pt,
               fill=gray!8, text width=5.4cm, align=left, inner sep=4pt,
               font=\footnotesize},
  final/.style={stage, fill=tealhead!18, line width=1.0pt},
  arr/.style={-{Latex[length=2mm]}, draw=tealrule, line width=0.7pt},
  sarr/.style={-{Latex[length=1.6mm]}, draw=gray!60, line width=0.5pt},
  node distance=7mm
]
\node[stage] (s1) {Records identified via database search\\[1pt]\textbf{$\sim$320}};
\node[stage, below=of s1] (s2) {After deduplication\\[1pt]\textbf{$\sim$240}};
\node[stage, below=of s2] (s3) {Title and abstract screened\\[1pt]\textbf{$\sim$150}};
\node[stage, below=of s3] (s4) {Full text assessed\\[1pt]\textbf{$\sim$120}};
\node[final, below=of s4] (s5) {\textbf{Included and cited: 105}};

\draw[arr] (s1) -- (s2);
\draw[arr] (s2) -- (s3);
\draw[arr] (s3) -- (s4);
\draw[arr] (s4) -- (s5);

\node[excl, right=30mm of $(s1)!0.5!(s2)$] (e1) {Duplicates across the four databases ($\sim$80)};
\node[excl, right=30mm of $(s2)!0.5!(s3)$] (e2) {Off topic on title or abstract ($\sim$90)};
\node[excl, right=30mm of $(s3)!0.5!(s4)$] (e3) {No evaluative use of a language model ($\sim$30)};
\node[excl, right=30mm of $(s4)!0.5!(s5)$] (e4) {Scope mismatch: reward modeling without judge evaluation, or $>$70B only ($\sim$15)};

\draw[sarr] ($(s1)!0.5!(s2)$) -- (e1);
\draw[sarr] ($(s2)!0.5!(s3)$) -- (e2);
\draw[sarr] ($(s3)!0.5!(s4)$) -- (e3);
\draw[sarr] ($(s4)!0.5!(s5)$) -- (e4);

\node[below=4mm of s5, text width=8.6cm, align=center, font=\footnotesize] {
  Of the 105 included: 29 judge systems, 20 benchmarks and metrics,
  26 training and methodology, 13 bias and robustness,
  17 foundational and model reports.};
\end{tikzpicture}
\end{adjustbox}

\caption{Literature selection flow. Left column gives records retained at each stage; right column gives the dominant exclusion reason. Only the final count is exact and can be checked against the reference list; upstream counts are recorded from the search logs and are approximate, since the same work was often retrieved by several query families.}
\label{fig:selection_flow}
\end{figure*}

\subsection{Search Strategy}

We searched Semantic Scholar, arXiv (cs.CL, cs.AI, cs.LG), ACL Anthology, and Google Scholar between November 2024 and July 2026 using the following query families:

\begin{enumerate}[leftmargin=1.5em, itemsep=2pt]
  \item \textbf{Core judging:} \texttt{("LLM-as-a-judge" OR "LLM judge" OR "language model evaluator" OR "automated evaluation")} combined with \texttt{("small" OR "7B" OR "8B" OR "3B" OR "13B" OR "14B" OR "efficient")}.
  \item \textbf{Training for evaluation:} \texttt{("reward model" OR "process reward" OR "GRPO" OR "DPO") AND ("judge" OR "evaluator" OR "scoring")}.
  \item \textbf{Multi-agent evaluation:} \texttt{("multi-agent debate" OR "ensemble judge" OR "panel of judges" OR "jury evaluation")}.
  \item \textbf{Reasoning budget:} \texttt{("overthinking" OR "token budget" OR "test-time compute" OR "chain-of-thought efficiency") AND ("evaluation" OR "judge")}.
  \item \textbf{Bias and robustness:} \texttt{("position bias" OR "self-preference" OR "adversarial" OR "prompt sensitivity") AND ("judge" OR "evaluator")}.
\end{enumerate}

We also followed citation chains from three anchor surveys \citep{gu2024survey, li2025generation, li2024llms} and tracked the RewardBench leaderboard~\citep{lambert2025rewardbench} for newly submitted judge systems.

\subsection{Inclusion and Exclusion Criteria}

\paragraph{Inclusion.} A paper was included if it met at least one of:
\begin{itemize}[leftmargin=1.2em, itemsep=2pt]
  \item It proposes, trains, or evaluates an SLM ($\leq$14B parameters) in an evaluative or scoring role.
  \item It introduces a benchmark, dataset, or metric used to evaluate language-model judges.
  \item It provides general-purpose methodology (multi-agent debate, ensemble voting, bias mitigation) that has been applied to or is directly applicable to SLM judges.
  \item It reports results for LLM judges ($>$14B) on benchmarks also used by SLM judge papers, enabling direct comparison. Such systems serve only as reference points and never count toward an insight's evidence grade.
\end{itemize}

\paragraph{Exclusion.} We excluded:
\begin{itemize}[leftmargin=1.2em, itemsep=2pt]
  \item Papers where language models are used purely as data annotators (e.g., labeling NER datasets) rather than as quality judges.
  \item Papers focused exclusively on reward modeling for RLHF without evaluation of the model as a standalone judge.
  \item Non-peer-reviewed blog posts, unless they describe a released system with reproducible results (e.g., model cards on HuggingFace with benchmark numbers).
  \item Preprints whose central claims could not be verified from the manuscript or released artifacts.
\end{itemize}

\subsection{Selection Accounting}

Screening followed a two-stage procedure for literature surveys. Records were first screened on title and abstract against the criteria above, and the remaining candidates were read in full before a final inclusion decision. Works found through citation chasing went through the same full-text check. Figure~\ref{fig:selection_flow} summarizes the selection process.

The largest source of exclusion at the full-text stage was scope mismatch: papers about reward modeling for RLHF that did not evaluate the model as a judge, and papers about LLM evaluation that studied only models above 70B without reporting results on smaller systems.

\subsection{Corpus Statistics}
\label{sec:corpus_stats}

We characterize the surveyed literature along four dimensions.

\paragraph{By training paradigm.} Table~\ref{tab:master_systems} groups the 19 judge systems into five paradigms: off-the-shelf prompting (1), supervised fine-tuning from distillation, human labels, or rubrics (8), self-supervised or preference optimization (1), reinforcement learning with verifiable rewards (7), and representation-based probing (2). Table~\ref{tab:training_overview} decomposes the same space more finely, separating the three supervised variants and listing preference optimization on its own. The RL systems all come from 2025--2026 work while the supervised systems cluster in 2023--2024, which is the decreasing-supervision gradient developed in Appendix~\ref{app:training}.

\paragraph{By base model family.} Qwen backbones (Qwen2.5, Qwen~3) support 8 of the 19 systems and LLaMA variants support 5, making these the two dominant bases. Mistral, Phi-3.5-mini, Vicuna, PaLM~2, and R1-Distill each support one. One system, J1, reports on both Qwen and LLaMA backbones and is counted under each, so the 18 family assignments cover 17 systems; the two representation-based systems probe a frozen model without committing to a family. The concentration on Qwen and LLaMA likely reflects permissive licensing and strong instruction-following baselines rather than any architectural advantage for judging.

\paragraph{By evaluation domain.} Nine of the 19 systems report on general-purpose or broad-coverage evaluation. Five target reasoning or verifiable tasks, and the remaining five each target a single setting: rubric-conditioned assessment, multilingual evaluation, mathematical process supervision, structured tasks, and fine-grained scoring. No system in the table is dedicated to safety, creative, or medical evaluation; those domains enter the survey only through benchmark and analysis papers (Appendix~\ref{app:domains}). The imbalance is itself a finding: evidence is dense where answers can be checked and thin everywhere else (Insight~I5).

\paragraph{By publication year.} Of the 105 cited works, 21 are from 2023 or earlier, 26 from 2024, 44 from 2025, and 14 from 2026. The subfield is young: every judge system in Table~\ref{tab:master_systems} carries a bibliography date of 2024 or later, and the two earliest of them, PandaLM and JudgeLM, first circulated as 2023 preprints.
